\documentclass[11pt,a4paper]{article}
\newcommand{\resultsdir}{../../results/task1}
% Common preamble of the task reports (compile with XeLaTeX).
\usepackage{fontspec}
\usepackage{polyglossia}
\setmainlanguage{polish}
\usepackage{amsmath,amssymb}
\usepackage[a4paper,margin=2.3cm]{geometry}
\usepackage{graphicx}
\usepackage{booktabs}
\usepackage{longtable}
\usepackage{array}
\usepackage{multirow}
\usepackage{adjustbox}
\usepackage{makecell}
\usepackage{placeins}
\usepackage{float}
\usepackage[font=small,labelfont=bf]{caption}
\usepackage{subcaption}
\usepackage{xcolor}
\usepackage{enumitem}
\usepackage[hidelinks]{hyperref}

% Without the texlive-lang-polish package XeLaTeX would hyphenate Polish
% words with English patterns, so hyphenation is switched off and looser
% inter-word spacing is allowed instead. With Polish patterns installed the
% next two penalties can be removed.
\hyphenpenalty=10000
\exhyphenpenalty=10000
\tolerance=3000
\emergencystretch=3em
\setlist{itemsep=2pt,topsep=4pt}
\setlength{\parskip}{3pt}
\renewcommand{\arraystretch}{1.12}

\newcommand{\tab}[1]{\input{\resultsdir/tables/#1}}
% table scaled down to the text width only when it is wider
\newcommand{\fittab}[1]{\adjustbox{max width=\textwidth}{\tab{#1}}}
\newcommand{\fig}[1]{\resultsdir/figures/#1}
\newcommand{\code}[1]{\texttt{#1}}
% left-aligned paragraph column (no stretched spaces in narrow cells)
\newcolumntype{L}[1]{>{\raggedright\arraybackslash}p{#1}}

% \maketitlepage{task title}{subtitle}
\newcommand{\maketitlepage}[2]{%
  \begin{titlepage}
    \centering
    {\large Politechnika Łódzka\\Wydział Fizyki Technicznej, Informatyki i Matematyki Stosowanej\par}
    \vspace{3.2cm}
    {\large Analiza danych złożonych z detekcją wyjątków\par}
    \vspace{0.8cm}
    {\LARGE\bfseries #1\par}
    \vspace{0.4cm}
    {\large #2\par}
    \vfill
    {\large Jędrzej Kostyk\par}
    {nr indeksu 264339\par}
    \vspace{1.2cm}
    {Rok akademicki 2026/2027\par}
  \end{titlepage}
}


\begin{document}
\maketitlepage{Zadanie 1}{Porównanie klasycznych metod detekcji wyjątków}
\tableofcontents
\newpage

\section{Cel i zakres}

Celem zadania jest porównanie pięciu klasycznych metod detekcji wyjątków
o różnych mechanizmach działania: Isolation Forest (izolacja przez losowe
podziały), Local Outlier Factor (lokalna gęstość), Connectivity-based Outlier
Factor (lokalna spójność sąsiedztwa), One-Class SVM (granica obszaru danych
w przestrzeni jądrowej) oraz DBSCAN (punkty szumu w klasteryzacji gęstościowej).
Szczególną uwagę poświęcono wpływowi hiperparametrów na liczbę i rodzaj
wykrywanych wyjątków, problemowi silnego niezbalansowania klas, różnicy
między wyjątkami globalnymi i lokalnymi oraz zgodności zbiorów obserwacji
wskazywanych przez poszczególne detektory.

Eksperymenty przeprowadzono na pięciu zbiorach z publicznych benchmarków
detekcji anomalii, dla 84 konfiguracji wymagających uczenia modelu
(154 po uwzględnieniu wariantów progu \code{contamination}) i łącznie
90 podziałów danych (7560 wytrenowanych modeli). Cały kod znajduje się
w katalogach \code{common/} i \code{task1\_classical/} repozytorium,
a wszystkie tabele i wykresy raportu są generowane skryptami
(sekcja~\ref{sec:repro}).

\section{Dane}

\subsection{Zbiory benchmarkowe}

Użyto gotowych zbiorów z trzech benchmarków detekcji anomalii:
ADBench~\cite{adbench}, ODDS~\cite{odds} i repozytorium DAMI~\cite{dami}.
Macierz cech $X$ i etykiety $y$ ($y=1$ -- anomalia) są wczytywane dokładnie
w postaci opublikowanej przez autorów benchmarków. \textbf{Etykiety anomalii
pochodzą z tych zbiorów, a nie z tego projektu}: w kodzie nie są generowane,
usuwane ani podpróbkowywane żadne obserwacje. Jedyne operacje na danych to
podział warstwowy i standaryzacja dopasowana na części uczącej
(sekcja~\ref{sec:protocol}). Zgodność plików sprawdzana jest sumami SHA256
(\code{data/CHECKSUMS.txt}) oraz udokumentowanymi kształtami i liczbą anomalii.

\begin{itemize}
  \item \textbf{Breastw} (ADBench, plik \code{4\_breastw.npz}) -- Breast Cancer
  Wisconsin (Original) z UCI~\cite{uci,breastw}: 683 biopsje opisane 9 cechami
  cytologicznymi w skali 1--10. Anomalie: guzy złośliwe (239, 35,0\,\%).
  \item \textbf{Mammography} (ADBench, \code{23\_mammography.npz}) -- dane
  z pracy Woods i in.~\cite{mammography} (zbiór spoza UCI, udostępniany
  w ODDS i ADBench): 11\,183 obszary mammogramów opisane 6 cechami obrazu.
  Anomalie: mikrozwapnienia (260, 2,3\,\%).
  \item \textbf{Wine} (ADBench, \code{45\_wine.npz}) -- Wine z
  UCI~\cite{uci,wine}: 13 wyników analizy chemicznej win. Autorzy ODDS
  przyjęli klasy 2 i 3 jako normalne (119), a klasę 1, zredukowaną przez
  nich do 10 obserwacji, jako anomalie (7,8\,\%).
  \item \textbf{Statlog Heart} (DAMI, plik
  \code{HeartDisease\_withoutdupl\_05\_v02.arff}) -- Statlog (Heart)
  z UCI~\cite{uci,heart}: 13 cech klinicznych. W wersji DAMI bez duplikatów,
  bez normalizacji, z 5\,\% anomalii (wariant v02): 150 osób zdrowych i 7 osób
  z chorobą serca jako anomalie (4,5\,\%). Normalizację wykonujemy sami na
  części uczącej.
  \item \textbf{Cardio} (ADBench, \code{6\_cardio.npz}) -- Cardiotocography
  z UCI~\cite{uci,cardio}: 1831 zapisów KTG opisanych 21 cechami. Klasa
  normalna (1655) jest normalna, klasa patologiczna zredukowana przez autorów
  ODDS do 176 obserwacji jest anomalią (9,6\,\%), klasa ,,podejrzana''
  została przez nich pominięta.
\end{itemize}

Zbiory różnią się wielkością (od 129 do 11\,183 obserwacji) i udziałem
anomalii (od 2,3 do 35\,\%). Breastw jest nietypowy: ponad jedna trzecia
obserwacji to anomalie, więc tworzą one dużą, zwartą grupę, a nie rzadkie
punkty. Wine i Statlog Heart mają tylko 10 i 7 anomalii, co ma konsekwencje
dla protokołu (sekcja~\ref{sec:protocol}).

\subsection{Wstępna analiza danych}

Tabela~\ref{tab:eda} podsumowuje zbiory. Żaden nie zawiera braków danych.
Cechy Breastw i Statlog Heart są całkowitoliczbowe (skale porządkowe, kody
kategorii), a pozostałe w większości ciągłe. Mammography i Cardio mają silnie
skośne rozkłady cech (średnia $|\gamma_1|$ odpowiednio 4,4 i 2,3;
rys.~\ref{fig:skew}). Wszystkie zbiory poza Statlog Heart zawierają silnie
skorelowane pary cech (maksymalna $|r|$ od 0,81 do 0,94), więc ich efektywna
wymiarowość jest mniejsza od liczby cech. Odsetek wartości $|z|>3$
(jednowymiarowe wartości odstające) wynosi 0,5--1,3\,\% i występują one
także w klasie normalnej, więc reguły typu ,,3 sigma'' nie wystarczą do
identyfikacji anomalii.

\begin{table}[H]
  \centering\small
  \caption{Charakterystyka zbiorów. Glob./lok. -- liczba anomalii globalnych
  i lokalnych wg definicji z sekcji~\ref{sec:kinds}; $\overline{|\gamma_1|}$ --
  średnia bezwzględna skośność cech; $|z|>3$ -- odsetek wartości odstających
  jednowymiarowo; max $|r|$ -- największa bezwzględna korelacja dwóch cech;
  PC 90\,\% -- liczba składowych głównych wyjaśniających 90\,\% wariancji.}
  \label{tab:eda}
  \resizebox{\textwidth}{!}{\tab{eda_summary}}
\end{table}

Rysunek~\ref{fig:eda_pca} przedstawia rzut danych na dwie pierwsze składowe
główne. W Breastw anomalie tworzą osobny, rozległy obszar. W Wine leżą na
skraju chmury danych normalnych jako niewielkie skupisko. W Mammography
i Cardio anomalie lokalne leżą na brzegu masy danych normalnych, a tylko
nieliczne są wyraźnie oddalone. W Statlog Heart anomalie są wymieszane
z danymi normalnymi. Potwierdza to rys.~\ref{fig:shift}: w Breastw wszystkie
cechy mają duży efekt ($|d|>0{,}8$), w Wine i Cardio -- kilka cech,
a w Statlog Heart i Mammography tylko pojedyncze cechy przekraczają tę granicę.

\begin{figure}[H]
  \centering
  \includegraphics[width=\textwidth]{\fig{eda_pca}}
  \caption{Rzut zbiorów na dwie pierwsze składowe główne (dane
  standaryzowane); kolor anomalii oznacza typ (globalne/lokalne).}
  \label{fig:eda_pca}
\end{figure}

\begin{figure}[H]
  \centering
  \includegraphics[width=\textwidth]{\fig{eda_feature_shift}}
  \caption{Bezwzględna wartość $d$ Cohena (standaryzowana różnica średnich
  anomalii i obserwacji normalnych) dla kolejnych cech; linia przerywana --
  umowna granica dużego efektu (0,8).}
  \label{fig:shift}
\end{figure}

\begin{figure}[H]
  \centering
  \includegraphics[width=0.7\textwidth]{\fig{eda_skewness}}
  \caption{Rozkład skośności cech w poszczególnych zbiorach.}
  \label{fig:skew}
\end{figure}

\section{Protokół eksperymentalny}
\label{sec:protocol}

\begin{enumerate}
  \item \textbf{Podział danych.} Dla każdego zbioru i ziarna $s$ dane dzielone
  są warstwowo na część uczącą, walidacyjną i testową w proporcji
  50/25/25\,\%. Te same ziarna dają te same podziały we wszystkich trzech
  zadaniach.
  \item \textbf{Liczba podziałów.} 10 podziałów (ziarna 0--9) dla Breastw,
  Mammography i Cardio oraz 30 podziałów (ziarna 0--29) dla Wine i Statlog
  Heart. W tych dwóch zbiorach część walidacyjna i testowa zawiera tylko
  2--3 anomalie (Wine: 3 w części testowej, Statlog Heart: 2). Podział
  warstwowy gwarantuje co najmniej jedną anomalię w każdej części, co jest
  dodatkowo sprawdzane w kodzie. Pojedynczy podział jest jednak bardzo
  zaszumiony, więc wynik uśredniono po trzykrotnie większej liczbie
  podziałów.
  \item \textbf{Przygotowanie.} Standaryzacja (\code{StandardScaler}) jest
  dopasowywana wyłącznie do części uczącej i stosowana do pozostałych części.
  \item \textbf{Uczenie.} Detektory uczone są \emph{bez etykiet} na części
  uczącej, która zawiera anomalie w naturalnej proporcji (scenariusz
  nienadzorowany). Każdy detektor zwraca wynik anomalii (wyższy = bardziej
  anomalny) dla nowych obserwacji.
  \item \textbf{Strojenie.} Etykiety części walidacyjnej służą do wyboru
  konfiguracji (najwyższa średnia PR-AUC na walidacji, uśredniona po
  podziałach) i progu decyzyjnego (maksymalizacja F1 na walidacji).
  Część testowa służy wyłącznie do końcowej oceny.
  \item \textbf{Metryki.} Precision, Recall, F1 (dla klasy ,,anomalia''),
  ROC-AUC, PR-AUC (\emph{average precision}), Accuracy (tylko do
  ilustracji jej wad) oraz czas uczenia i czas oceny części testowej.
  Raportowane są średnie (i odchylenia standardowe) po podziałach.
\end{enumerate}

Dla metod z parametrem \code{contamination} (IForest, LOF, COF) parametr
ten nie zmienia wyniku anomalii, a jedynie próg: oznaczane jest
$c\cdot 100\,\%$ obserwacji uczących o najwyższym wyniku. Dlatego wszystkie
wartości $c$ oceniane są na jednym wytrenowanym modelu. W przypadku OCSVM
decyzja wynika z granicy SVM (sterowanej przez $\nu$), a w DBSCAN -- z
promienia $\varepsilon$. W tabelach siatki (dodatek~\ref{app:grid}) Precision,
Recall i F1 dotyczą tych ,,wbudowanych'' reguł decyzyjnych, a w porównaniu
najlepszych konfiguracji (sekcja~\ref{sec:best}) -- progu wybranego na walidacji.

\section{Metody i badane hiperparametry}

\paragraph{Isolation Forest.} Las $T$ drzew dzielących dane losowymi cięciami
losowych cech; wynik anomalii zależy od średniej głębokości, na której
obserwacja zostaje odizolowana. Punkty odległe od masy danych izolowane są
po niewielu cięciach. Metoda patrzy na dane \emph{globalnie} -- cięcia
osiowe nie ,,widzą'' lokalnych różnic gęstości. Badano
$T\in\{50,100,200,500\}$ i $c\in\{0{,}02;\,0{,}05;\,0{,}1;\,0{,}15;\,0{,}2;\,0{,}35\}$
(wartość 0,35 odpowiada udziałowi anomalii w Breastw).

\paragraph{LOF.} Stosunek lokalnej gęstości osiągalności sąsiadów punktu do
jego własnej gęstości. Wartości $\gg 1$ oznaczają punkt rzadszy niż jego
otoczenie -- miara jest \emph{względna}, więc wykrywa wyjątki lokalne przy
gęstych skupiskach. Badano $k\in\{5,10,20,35,50\}$ i te same wartości $c$.
Użyto trybu \code{novelty=True}, który pozwala oceniać nowe dane.

\paragraph{COF.} Wariant LOF, w którym gęstość zastąpiono \emph{spójnością}:
średnia odległość łańcuchowa (\emph{average chaining distance}) wzdłuż
ścieżki najbliższych połączeń (kolejność algorytmu Prima) w zbiorze $k$
sąsiadów. COF lepiej niż LOF radzi sobie z wzorcami o niskiej gęstości,
ale regularnej strukturze (np. punkty ułożone wzdłuż linii). Implementacja
została napisana od podstaw (\code{common/cof.py}), bo dostępne
implementacje nie oceniają nowych danych; zgodność z definicją sprawdzono
na naiwnej implementacji referencyjnej (różnica 0). Siatka jak dla LOF.

\paragraph{One-Class SVM.} Wyznacza w przestrzeni cech jądra hiperpłaszczyznę
oddzielającą dane od początku układu z maksymalnym marginesem; $\nu$ jest
górnym ograniczeniem odsetka punktów uczących poza granicą (i dolnym
ograniczeniem odsetka wektorów nośnych). Badano jądra RBF, wielomianowe
(stopnia 3), sigmoidalne i liniowe, $\nu\in\{0{,}01;\,0{,}05;\,0{,}1;\,0{,}2;\,0{,}5\}$
oraz $\gamma\in\{0{,}1/d,\ 1/d,\ 10/d\}$ (50 konfiguracji).

\paragraph{DBSCAN.} Punkty rdzeniowe mają co najmniej minPts sąsiadów
w promieniu $\varepsilon$; punkty nieosiągalne z żadnego rdzenia są szumem
(anomalią). Aby oceniać nowe dane, wynikiem anomalii jest odległość do
najbliższego punktu rdzeniowego z części uczącej, a anomalią -- punkt, dla
którego odległość ta przekracza $\varepsilon$. Ponieważ skala odległości
zależy od wymiaru i zbioru, $\varepsilon$ ustalano jako kwantyl
$q_\varepsilon\in\{0{,}5;\,0{,}75;\,0{,}9;\,0{,}95;\,0{,}99\}$ rozkładu
odległości do minPts-tego sąsiada (heurystyka wykresu $k$-odległości),
minPts $\in\{3,5,10,20\}$.

\section{Wyniki}

\subsection{Dlaczego Accuracy prowadzi do błędnych wniosków}

Tabela~\ref{tab:acc} zestawia trywialny klasyfikator ,,zawsze normalny''
z najlepszym i najsłabszym (wg F1) detektorem. Poza Breastw klasyfikator
trywialny osiąga Accuracy 0,90--0,98, nie wykrywając \emph{żadnej} anomalii
(Recall = F1 = 0). Na Mammography i Statlog Heart jego Accuracy (0,977
i 0,955) jest \emph{wyższa} niż najlepszego detektora (0,974 i 0,889),
choć ten wykrywa odpowiednio 23 i 48\,\% anomalii. Na Mammography różnica
Accuracy między najlepszym a najsłabszym detektorem to 0,005, a w F1 --
0,09. Przy udziale anomalii $\pi$ Accuracy jest zdominowana przez klasę
normalną: wykrycie wszystkich anomalii zmienia ją o co najwyżej $\pi$.
Dlatego w dalszej analizie podstawą są PR-AUC (której wartość dla losowego
rankingu wynosi $\pi$, a nie 0,5 jak dla ROC-AUC) oraz F1.

\begin{table}[H]
  \centering\small
  \caption{Accuracy a jakość wykrywania anomalii (część testowa, średnie
  po podziałach). ,,Losowy'' -- PR-AUC rankingu losowego (= udział anomalii).}
  \label{tab:acc}
  \resizebox{\textwidth}{!}{\tab{accuracy_trap}}
\end{table}

\subsection{Wpływ hiperparametrów}

\paragraph{Isolation Forest -- liczba drzew.} Liczba drzew praktycznie nie
wpływa na średnią jakość (rys.~\ref{fig:iforest}; zmiany PR-AUC $\le 0{,}045$,
największe na małych Wine i Statlog Heart, gdzie odchylenie między podziałami
wynosi 0,15--0,25). Na Mammography większa liczba drzew zmniejsza wariancję między
podziałami (odchylenie 0,066 dla 50 drzew i 0,045 dla 500). Ponieważ czas
uczenia rośnie liniowo z $T$, 100--200 drzew jest rozsądnym kompromisem.
Różnice między zbiorami są o rząd wielkości większe niż różnice wynikające
z $T$, co pokazuje, że ograniczeniem IForest jest jego mechanizm, a nie
liczba drzew.

\begin{figure}[H]
  \centering
  \includegraphics[width=0.85\textwidth]{\fig{iforest_estimators}}
  \caption{Isolation Forest: średnia i rozrzut PR-AUC w zależności od liczby
  drzew.}
  \label{fig:iforest}
\end{figure}

\paragraph{LOF i COF -- liczba sąsiadów.} Wpływ $k$ jest silny i zależy od
struktury anomalii (rys.~\ref{fig:knn}). W Wine PR-AUC LOF rośnie z 0,10
($k=5$) do 0,49 ($k=20$): anomalie tworzą małe skupisko (ok. 5 punktów
w części uczącej), więc przy $k=5$ sąsiadami anomalii są inne anomalie i jej
gęstość względna wydaje się normalna (efekt maskowania). Dopiero $k$ większe
niż rozmiar skupiska pozwala je wykryć. W Cardio, Mammography i Breastw
jakość rośnie monotonicznie aż do $k=50$ (Cardio: 0,20 $\rightarrow$ 0,31),
bo anomalie tworzą tam większe grupy. Skrajnym przypadkiem jest Breastw:
239 anomalii tworzy własny gęsty obszar, a LOF i COF niezależnie od $k$
osiągają PR-AUC tylko 0,29--0,43 przy bazowym 0,35. Guz złośliwy jest
,,normalny'' względem innych guzów złośliwych, a metody lokalne porównują
punkt właśnie z jego sąsiadami. COF jest prawie wszędzie słabszy od LOF
(np. Cardio: 0,20 wobec 0,31) -- ścieżka łańcuchowa przy dużym $k$
,,przechodzi'' przez sąsiednie skupiska i zaciera strukturę.

\begin{figure}[H]
  \centering
  \includegraphics[width=\textwidth]{\fig{neighbors_pr_auc}}
  \caption{LOF i COF: PR-AUC (część testowa) w funkcji liczby sąsiadów;
  pasy -- $\pm 1$ odchylenie standardowe między podziałami.}
  \label{fig:knn}
\end{figure}

\paragraph{Parametr \code{contamination}.} Parametr ten ustala odsetek
oznaczanych obserwacji, a więc bezpośrednio steruje kompromisem
Precision--Recall (rys.~\ref{fig:cont}). Dla IForest na Cardio wzrost $c$
z 0,02 do 0,35 podnosi Recall z 0,17 do 0,995, a Precision spada z 0,84 do
0,26. F1 jest największe w pobliżu rzeczywistego udziału anomalii: na Cardio
($\pi=0{,}096$) przy $c=0{,}15$ (0,57), na Breastw ($\pi=0{,}35$) przy
$c=0{,}35$ (0,93), na Mammography ($\pi=0{,}023$) przy $c=0{,}02$ (0,25).
Ważna obserwacja: contamination \emph{nie zmienia} rankingu obserwacji
(ROC-AUC i PR-AUC są stałe) -- zmienia tylko, ile najwyżej ocenionych punktów
nazywamy anomaliami. Przyjęcie $c$ mniejszego od rzeczywistego udziału
daje wysoką precyzję i niski Recall (detektor wskazuje tylko najbardziej
oczywiste, zwykle globalne wyjątki), a większego -- odwrotnie. W praktyce
$c$ jest nieznane, dlatego w porównaniu metod próg wyznaczano na walidacji.

\begin{figure}[H]
  \centering
  \includegraphics[width=\textwidth]{\fig{contamination_effect}}
  \caption{Precision i Recall (część testowa) w zależności od contamination
  dla najlepszych (wg walidacji) wartości $T$ i $k$; linia przerywana --
  rzeczywisty udział anomalii.}
  \label{fig:cont}
\end{figure}

\paragraph{One-Class SVM.} Kluczowe okazały się jądro i $\gamma$
(rys.~\ref{fig:ocsvm_heat} i~\ref{fig:ocsvm_kern}). Przy $\nu$ i $\gamma$
wybranych na walidacji jądro RBF jest najlepsze na Breastw (PR-AUC 0,94)
i Statlog Heart (0,42), a sigmoidalne -- na Wine (0,89), Cardio (0,76)
i Mammography (0,22). Jądro liniowe (0,13--0,59) i wielomianowe (0,02--0,46)
są słabe -- liniowa granica oddzielająca dane od początku układu nie może
opisać skupisk ani obszarów niewypukłych, a jądro wielomianowe mocno zależy
od skali danych. Sukces jądra sigmoidalnego jest empiryczny: nie jest ono
dodatnio określone, więc wynik nie ma interpretacji granicy obszaru
w przestrzeni Hilberta. Dla danych standaryzowanych przy $\nu=0{,}5$
(we wszystkich trzech przypadkach) wynik zależy od iloczynów skalarnych
z wieloma wektorami nośnymi i działa jak miara odchylenia od dominujących
kierunków danych. Dla RBF $\gamma$ steruje ,,lokalnością'' granicy:
$\gamma=10/d$ daje granicę ciasno opinającą dane i jest najlepsze dla
Breastw (0,92 niezależnie od $\nu$), a $\gamma=0{,}1/d$ -- gładką granicę,
najlepszą dla Cardio (0,60 przy $\nu=0{,}5$, wobec 0,26 przy $\gamma=10/d$).

Parametr $\nu$ ogranicza odsetek \emph{uczących} punktów poza granicą, ale
na nowych danych odsetek oznaczonych punktów może być znacznie większy
(rys.~\ref{fig:ocsvm_nu}). Na małych Wine i Statlog Heart (64 i 78 punktów
uczących, 13 cech) przy $\gamma=1/d$ OCSVM oznacza 41--46\,\% punktów
testowych nawet dla $\nu=0{,}01$. Granica jest tak ciasno dopasowana do
punktów uczących, że nowe punkty często leżą poza nią (nadmierne
dopasowanie). Na większych Cardio i Mammography odsetek ten jest bliski
$\nu$ dopiero od $\nu\approx 0{,}1$. Przy $\gamma=10/d$ wynik prawie nie
zależy od $\nu$: jądro jest tak wąskie, że niemal każdy punkt staje się
wektorem nośnym, a funkcja decyzyjna OCSVM sprowadza się do jądrowego
estymatora gęstości (Parzena).

\begin{figure}[H]
  \centering
  \includegraphics[width=\textwidth]{\fig{ocsvm_rbf_heatmap}}
  \caption{OCSVM z jądrem RBF: PR-AUC (część testowa) w zależności od
  $\nu$ i $\gamma$.}
  \label{fig:ocsvm_heat}
\end{figure}

\begin{figure}[H]
  \centering
  \includegraphics[width=0.8\textwidth]{\fig{ocsvm_kernels}}
  \caption{OCSVM: PR-AUC (część testowa) dla każdego jądra przy $\nu$
  i $\gamma$ wybranych na walidacji.}
  \label{fig:ocsvm_kern}
\end{figure}

\begin{figure}[H]
  \centering
  \includegraphics[width=0.85\textwidth]{\fig{ocsvm_nu}}
  \caption{OCSVM (RBF, $\gamma=1/d$): odsetek oznaczonych punktów testowych
  w funkcji $\nu$ (linia przerywana -- $y=\nu$) oraz kompromis
  Precision--Recall przy zmianie $\nu$.}
  \label{fig:ocsvm_nu}
\end{figure}

\paragraph{DBSCAN.} Rysunek~\ref{fig:dbscan} pokazuje, że w DBSCAN o
wynikach decyduje przede wszystkim $\varepsilon$: odsetek punktów szumu
maleje od 29--46\,\% dla $q_\varepsilon=0{,}5$ do 0,4--1,1\,\% dla
$q_\varepsilon=0{,}99$. F1 osiąga maksimum, gdy odsetek szumu jest bliski
udziałowi anomalii lub nieco większy: Breastw $q_\varepsilon=0{,}5$
(szum 46\,\%, F1 = 0,89), Cardio $q_\varepsilon=0{,}75$ (18\,\%, 0,52),
Mammography $q_\varepsilon=0{,}95$ (3\,\%, 0,26). Wpływ minPts najlepiej
widać na Wine. Anomalie (ok. 5 punktów w części uczącej) tworzą małe
skupisko, które przy minPts $\le 5$ staje się osobnym klastrem: punkty
rdzeniowe są wtedy wśród anomalii i dla $q_\varepsilon\ge 0{,}75$ F1 spada
do zera. Dopiero minPts = 10 lub 20 (więcej niż rozmiar skupiska) sprawia,
że anomalie nie mogą utworzyć rdzenia i stają się szumem (F1 = 0,53--0,54
przy $q_\varepsilon=0{,}5$). Wadą DBSCAN jest jeden globalny promień
$\varepsilon$ -- nie da się dobrać go jednocześnie do gęstego i rzadkiego
skupiska.

\begin{figure}[H]
  \centering
  \includegraphics[width=\textwidth]{\fig{dbscan_heatmap}}
  \caption{DBSCAN: F1 (górny wiersz) i odsetek punktów testowych uznanych
  za szum (dolny wiersz) w zależności od minPts i kwantyla $q_\varepsilon$
  wyznaczającego $\varepsilon$.}
  \label{fig:dbscan}
\end{figure}

\subsection{Porównanie najlepszych konfiguracji}
\label{sec:best}

Tabela~\ref{tab:best} i rys.~\ref{fig:best} przedstawiają wyniki konfiguracji
wybranych na walidacji, z progiem maksymalizującym F1 na walidacji.

\begin{table}[H]
  \centering\small
  \caption{Najlepsze konfiguracje (wybór na walidacji wg PR-AUC, próg
  z walidacji wg F1); część testowa, średnia ($\pm$ odchylenie std.)
  po podziałach (10 lub 30). Czas = uczenie + ocena części testowej.
  Pogrubienie -- najlepsza metoda na zbiorze.}
  \label{tab:best}
  \resizebox{\textwidth}{!}{\tab{best_results}}
\end{table}

\begin{figure}[H]
  \centering
  \includegraphics[width=\textwidth]{\fig{best_comparison}}
  \caption{PR-AUC i F1 najlepszych konfiguracji; słupki błędów --
  odchylenie standardowe między podziałami.}
  \label{fig:best}
\end{figure}

Nie ma metody najlepszej na wszystkich zbiorach. OCSVM wygrywa na Cardio
(PR-AUC 0,76, F1 0,73) i Wine (0,89; 0,77). DBSCAN wygrywa na Breastw
(0,99; 0,94) i ma najwyższe PR-AUC i F1 na Mammography (0,23; 0,29),
ale tam różnice między DBSCAN, OCSVM i IForest (0,21--0,23) mieszczą się
w odchyleniu między podziałami. Na Statlog Heart wszystkie metody osiągają
PR-AUC 0,39--0,44 przy odchyleniu 0,22--0,25, więc nie da się ich
uporządkować. Metody lokalne (LOF, COF) są wyraźnie najsłabsze na Breastw
i Cardio. Isolation Forest jest solidny, ale nigdzie nie wygrywa.

Na tle wysokich ROC-AUC widać, dlaczego ROC-AUC jest metryką optymistyczną
przy niezbalansowaniu. Na Mammography wszystkie metody mają ROC-AUC
0,74--0,85, a PR-AUC tylko 0,11--0,23: IForest z ROC-AUC 0,85 osiąga
F1 = 0,24, czyli przy progu z walidacji trzy na cztery zgłoszenia są
fałszywe. Oś FPR na krzywej ROC jest normalizowana liczbą obserwacji
normalnych (42 razy większą od liczby anomalii), więc dziesiątki fałszywych
alarmów prawie nie obniżają ROC-AUC. PR-AUC i Precision bezpośrednio
pokazują ten koszt. Na Statlog Heart ROC-AUC $\approx 0{,}8$ odpowiada
PR-AUC $\approx 0{,}4$ przy bazowym 0,045 -- ranking jest znacznie lepszy od
losowego, ale część testowa zawiera tylko dwie anomalie, więc F1 przyjmuje
w pojedynczym podziale kilka możliwych wartości (odchylenie 0,20--0,26).

Rozbieżność między siatką a wyborem na walidacji jest niewielka: dla 18
z 25 par zbiór--metoda konfiguracja wybrana na walidacji jest też najlepsza
na teście, a w pozostałych różnica PR-AUC nie przekracza 0,02 -- z wyjątkiem
DBSCAN na Statlog Heart (najlepsza na teście 0,50, wybrana na walidacji 0,42).
Część walidacyjna zawiera tam tylko 2 anomalie, więc wybór hiperparametrów
jest zaszumiony. Raportowanie maksimum z siatki na zbiorze testowym
zawyżałoby wyniki -- dlatego wszystkie porównania metod opierają się na
wyborze walidacyjnym.

\subsection{Wyjątki lokalne i globalne}
\label{sec:kinds}

Benchmarki nie określają typu anomalii, więc przyjęto definicję operacyjną:
anomalia jest \emph{globalna}, jeśli jej odporna odległość Mahalanobisa od
klasy normalnej przekracza 99. percentyl odległości obserwacji normalnych,
a w przeciwnym razie \emph{lokalna} (leży wewnątrz globalnej ,,otoczki''
danych normalnych i odróżnia ją tylko sąsiedztwo). Kowariancję klasy
normalnej szacowano estymatorem MCD. W Cardio, które ma liniowo zależne cechy
(macierz kowariancji źle uwarunkowana), użyto estymatora Ledoita--Wolfa.
Etykiety typów służą wyłącznie analizie -- nie zmieniają etykiet benchmarku.
Większość anomalii jest w tym sensie lokalna (tabela~\ref{tab:eda}):
93\,\% w Mammography, 90\,\% w Wine, 86\,\% w Statlog Heart i 65\,\%
w Cardio. Wyjątkiem jest Breastw (34\,\% lokalnych).

Rysunek~\ref{fig:regions} pokazuje obszary decyzyjne w dwuwymiarowym rzucie
części uczącej Cardio na dwie pierwsze składowe główne. Wszystkie metody
oznaczają 10\,\% punktów, co odpowiada udziałowi anomalii w Cardio.
IForest, OCSVM i DBSCAN wyznaczają \emph{jeden} zwarty obszar normalności
wokół masy danych. LOF i COF tworzą postrzępioną granicę, która oznacza
także luki wewnątrz chmury i obrzeża mniejszych zgęszczeń. Rzut zachowuje
jednak tylko część struktury 21-wymiarowych danych: wiele anomalii
lokalnych leży w nim wewnątrz chmury danych normalnych, a nawet najlepszy
w tym rzucie IForest osiąga F1 = 0,43. Rysunek ilustruje więc kształt
granic poszczególnych metod, a nie ich skuteczność w pełnej przestrzeni.

\begin{figure}[H]
  \centering
  \includegraphics[width=\textwidth]{\fig{decision_regions_2d}}
  \caption{Obszary uznane za anomalne (kolor jasnopomarańczowy) w rzucie
  części uczącej Cardio (podział 0) na dwie pierwsze składowe główne;
  detektory uczone w tym rzucie. Kolor punktów -- typ anomalii wyznaczony
  w pełnej przestrzeni cech. R$_{lok}$, R$_{glob}$ -- Recall anomalii
  lokalnych i globalnych.}
  \label{fig:regions}
\end{figure}

W pełnej przestrzeni (tabela~\ref{tab:kinds}, rys.~\ref{fig:kinds})
anomalie globalne są wykrywane prawie w całości przez wszystkie metody
(Recall 0,89--1,00; jedynym wyjątkiem jest COF na Cardio, 0,59).
O jakości decydują anomalie lokalne. Na Cardio najlepiej wykrywa je LOF
(Recall 0,79), wyraźnie przed OCSVM (0,66) i IForest (0,60), a DBSCAN i COF
wykrywają je słabo (0,31 i 0,30). Na Statlog Heart LOF również prowadzi
(0,39). Na Wine i Mammography najwięcej anomalii lokalnych wykrywa OCSVM
(0,74 i 0,27). W Mammography wszystkie metody wykrywają mniej niż 30\,\%
anomalii lokalnych, które stanowią 93\,\% anomalii tego zbioru. To wyjaśnia
niskie PR-AUC wszystkich metod na Mammography. Na Breastw podział na
globalne i lokalne traci znaczenie: progi są tak niskie (35\,\% anomalii),
że oba typy są wykrywane w 87--100\,\%. W Wine i Statlog Heart anomalii
globalnych jest tylko 6 i 9 (suma po wszystkich częściach testowych), więc
dla tych zbiorów Recall globalnych należy traktować orientacyjnie.

\begin{table}[H]
  \centering\small
  \caption{Recall według typu anomalii (część testowa; $n$ -- liczba anomalii
  danego typu zsumowana po częściach testowych wszystkich podziałów).}
  \label{tab:kinds}
  \tab{recall_by_kind}
\end{table}

\begin{figure}[H]
  \centering
  \includegraphics[width=\textwidth]{\fig{recall_by_kind}}
  \caption{Recall według typu anomalii dla najlepszych konfiguracji.}
  \label{fig:kinds}
\end{figure}

\subsection{Zgodność zbiorów wykrytych wyjątków}

Zgodność zbiorów obserwacji oznaczonych przez detektory zmierzono
indeksem Jaccarda $J(A,B)=|A\cap B|/|A\cup B|$ (rys.~\ref{fig:jac}, średnia
po podziałach). Średnia zgodność par wynosi 0,62 w Breastw, ale tylko
0,28--0,38 w pozostałych zbiorach -- gdy problem jest trudny, detektory
wskazują różne obserwacje. Wyraźnie widać dwie rodziny metod. Metody
globalne IForest, OCSVM i DBSCAN są ze sobą zgodne (Breastw 0,89--0,92,
Cardio 0,50--0,55), a lokalne LOF i COF -- ze sobą (Breastw 0,78,
Mammography 0,49). Między rodzinami zgodność jest niska (Breastw 0,42--0,48,
Cardio 0,14--0,43). Zgodność odzwierciedla więc mechanizm metody bardziej
niż jej jakość.

\begin{figure}[H]
  \centering
  \includegraphics[width=\textwidth]{\fig{jaccard_heatmap}}
  \caption{Indeks Jaccarda zbiorów obserwacji testowych oznaczonych jako
  anomalie przez pary detektorów (średnia po podziałach).}
  \label{fig:jac}
\end{figure}

Rysunek~\ref{fig:overlap} pokazuje, iloma detektorami wskazywana jest
każda prawdziwa anomalia i każdy fałszywy alarm. W Breastw 83\,\% anomalii
wskazują wszystkie pięć metod i żadna nie zostaje pominięta przez
wszystkie. W pozostałych zbiorach zgodnie wskazywanych jest tylko
10--27\,\% anomalii, a żaden detektor nie wykrywa 3\,\% anomalii Cardio,
7\,\% Wine, 43\,\% Statlog Heart i 58\,\% Mammography. Tych obserwacji
nie uratuje żadna kombinacja detektorów. Fałszywe alarmy są natomiast
w dużej mierze \emph{indywidualne}: 42--59\,\% fałszywych alarmów zgłasza
tylko jeden detektor (poza Breastw, 31\,\%). Anomalie wykrywane wyłącznie
przez jedną metodę to głównie trafienia LOF (30 przypadków w Cardio)
i OCSVM (39 w Mammography, 14 w Cardio, 11 w Wine). Ta asymetria -- błędy
indywidualne, część trafień unikatowa -- jest przesłanką, że łączenie
detektorów może poprawić jakość; zostanie to zbadane w zadaniu 3.

\begin{figure}[H]
  \centering
  \includegraphics[width=\textwidth]{\fig{detector_overlap}}
  \caption{Rozkład liczby detektorów (spośród 5) wskazujących prawdziwe
  anomalie (lewo) oraz fałszywe alarmy wskazane przez co najmniej jeden
  detektor (prawo); część testowa, wszystkie podziały.}
  \label{fig:overlap}
\end{figure}

\subsection{Kompromis jakość -- czas}

Rysunek~\ref{fig:time} zestawia czas (uczenie + ocena części testowej)
z PR-AUC dla wszystkich konfiguracji. Na małych zbiorach (Wine, Statlog
Heart; ok. 70 punktów uczących) OCSVM i LOF potrzebują 1--2~ms, a Isolation
Forest ok. 145~ms. Jego koszt zależy głównie od liczby drzew, a każde drzewo
budowane jest na podpróbce 256 punktów, więc prawie nie zależy od $n$.
Wraz z rozmiarem danych proporcje się odwracają. Mediana czasu OCSVM rośnie
z 3~ms (Breastw, 341 punktów uczących) przez 17~ms (Cardio, 915) do 246~ms
(Mammography, 5591; maksimum 1,7~s), czyli szybciej niż liniowo, zgodnie
ze złożonością $O(n^2)$--$O(n^3)$ rozwiązania problemu kwadratowego.
IForest w tym samym zakresie zmienia się tylko z 158 do 214~ms. COF jest
najwolniejszy na większych zbiorach (mediana 343~ms na Mammography), bo dla
każdego punktu buduje ścieżkę łańcuchową w czasie $O(k^2)$. LOF (120~ms)
i DBSCAN (252~ms) wymagają wyszukiwania sąsiadów. W badanym zakresie
najkorzystniejszy stosunek jakości do czasu ma OCSVM na zbiorach
do kilku tysięcy obserwacji. Przy dziesiątkach tysięcy obserwacji przewagę
zyskuje IForest (por. pomiary dla Campaign w zadaniu 2).

\begin{figure}[H]
  \centering
  \includegraphics[width=\textwidth]{\fig{time_vs_quality}}
  \caption{Czas uczenia i oceny a PR-AUC (część testowa) dla wszystkich
  konfiguracji siatki (średnie po podziałach).}
  \label{fig:time}
\end{figure}

\section{Dyskusja}

\paragraph{Dlaczego algorytmy zachowują się inaczej?} Różnice wynikają
z tego, \emph{co} każda metoda uznaje za nietypowe.
\begin{itemize}
  \item \textbf{Isolation Forest} mierzy łatwość izolacji osiowymi cięciami
  w skali całego zbioru. Jest skuteczny, gdy anomalie odstają globalnie lub
  w pojedynczych cechach (Breastw: 0,97), i stabilny względem
  hiperparametrów, ale słabszy dla anomalii lokalnych (Cardio: Recall
  lokalnych 0,60 wobec 0,79 dla LOF).
  \item \textbf{LOF} porównuje gęstość punktu z gęstością sąsiadów, więc
  najlepiej wykrywa wyjątki lokalne (Cardio, Statlog Heart). Zawodzi, gdy
  anomalii jest dużo i tworzą własny gęsty obszar (Breastw: 0,43) -- wtedy
  założenie ,,anomalia jest rzadsza od sąsiadów'' nie jest spełnione.
  Za małe $k$ prowadzi do maskowania małych skupisk anomalii (Wine).
  \item \textbf{COF} zastępuje gęstość spójnością ścieżki; w danych
  bez wyraźnych struktur liniowych daje wyniki gorsze niż LOF, a jest
  najwolniejszy.
  \item \textbf{OCSVM} wyznacza jedną granicę dla całego zbioru; przy
  dobrze dobranym jądrze i $\gamma$ jest najlepszy na Cardio i Wine, ale
  wynik silnie zależy od jądra i $\gamma$, a na małych zbiorach granica
  nadmiernie dopasowuje się do punktów uczących.
  \item \textbf{DBSCAN} daje wysoką precyzję, gdy anomalie są wyraźnie
  oddzielone (Breastw: 0,93), bo za szum uznaje tylko punkty rzeczywiście
  odizolowane. Jeden globalny promień nie pasuje jednak jednocześnie do
  skupisk o różnej gęstości (Recall lokalnych w Cardio 0,31), a minPts musi
  przekraczać rozmiar skupisk anomalii.
\end{itemize}

\paragraph{Wpływ parametrów na rodzaj wykrywanych wyjątków.} Parametry
progowe (contamination, $\nu$, $\varepsilon$) decydują o \emph{liczbie}
wykrywanych wyjątków i przesuwają kompromis Precision--Recall. Ich
zmniejszanie sprawia, że detektor zgłasza tylko najbardziej oczywiste,
zwykle globalne wyjątki. Parametry kształtujące wynik ($k$ w LOF/COF,
$\gamma$ i jądro w OCSVM, minPts w DBSCAN) decydują o \emph{rodzaju}
wykrywanych wyjątków -- o skali, na której definiowana jest ,,normalność''.
Małe $k$ lub duże $\gamma$ czynią detektor lokalnym (wykrywa odchylenia
od najbliższego otoczenia, ale też maskuje małe skupiska anomalii), a duże
$k$ lub małe $\gamma$ -- globalnym.

\paragraph{Interpretacja metryk przy niezbalansowaniu.} Przy 2--10\,\%
anomalii Accuracy jest bezużyteczna (tabela~\ref{tab:acc}), a ROC-AUC
optymistyczna, bo fałszywe alarmy są normalizowane dużą liczbą obserwacji
normalnych (Mammography: ROC-AUC 0,85 przy F1 0,24). PR-AUC odnosi się do
bazowego poziomu $\pi$ i bezpośrednio uwzględnia koszt fałszywych alarmów,
dlatego użyto jej do wyboru konfiguracji. Precision i Recall zależą od
progu, więc porównując metody trzeba ustalić wspólną regułę wyboru progu
(tu: maksymalizacja F1 na walidacji). F1 traktuje oba rodzaje błędów
symetrycznie; w zastosowaniach (np. diagnostyka) należałoby dobrać próg do
kosztu pominięcia anomalii.

\paragraph{Ograniczenia.}
\begin{itemize}
  \item Małe liczby anomalii: części testowe Wine i Statlog Heart zawierają
  3 i 2 anomalie, części walidacyjne -- 2--3. Wybór hiperparametrów i progu
  jest więc obarczony dużą wariancją (odchylenie F1 0,19--0,26). Trzykrotnie
  większa liczba podziałów (30) zmniejsza niepewność średniej, ale nie
  zmniejsza zmienności pojedynczego wyniku; wniosków o kolejności metod na
  tych zbiorach nie należy wyciągać.
  \item Etykiety benchmarków są przyjęte bez zmian. Klasa uznana przez autorów
  benchmarku za anomalię nie musi być ,,wyjątkiem'' w sensie rozkładu
  (Statlog Heart, Mammography). Ocena mierzy wtedy zgodność anomalii
  statystycznych z tą klasą, a nie jakość detektora w ogóle. Breastw z 35\,\%
  anomalii jest raczej problemem dwóch skupisk niż detekcji rzadkich zdarzeń.
  \item Podział anomalii na globalne i lokalne opiera się na przyjętej
  definicji (MCD lub Ledoit--Wolf, 99. percentyl); inna definicja zmieniłaby
  proporcje typów.
  \item Siatki hiperparametrów są skończone; DBSCAN był parametryzowany
  kwantylem $\varepsilon$, a OCSVM -- względnym $\gamma$.
\end{itemize}

\paragraph{Kierunki dalszych badań.} Analiza wpływu reprezentacji danych
(zadanie 2), łączenie detektorów, które -- jak pokazała analiza zgodności --
należą do dwóch rodzin i popełniają w dużej mierze różne błędy (zadanie 3),
oraz metody półnadzorowane wykorzystujące niewielką liczbę etykiet
z walidacji.

\section{Wnioski}

\begin{enumerate}
  \item Żadna metoda nie dominuje: OCSVM jest najlepszy na Cardio (PR-AUC
  0,76) i Wine (0,89), DBSCAN na Breastw (0,99) i Mammography (0,23), a na
  Statlog Heart metody są nierozróżnialne. LOF i COF są wyraźnie najsłabsze
  tam, gdzie anomalie tworzą duże skupiska (Breastw).
  \item O tym, \emph{jakie} wyjątki są wykrywane, decyduje mechanizm metody:
  anomalie globalne wykrywają prawie w całości wszystkie metody, a różnice
  dotyczą anomalii lokalnych. Tu LOF prowadzi na Cardio (Recall 0,79 wobec
  0,30--0,66), a DBSCAN i COF są najsłabsze. Metody dzielą się na dwie
  rodziny (globalną i lokalną), co widać w zgodności ich decyzji.
  \item Najważniejsze hiperparametry to $k$ (LOF/COF), $\gamma$ i jądro (OCSVM)
  oraz $\varepsilon$ i minPts (DBSCAN); liczba drzew w IForest ma znikomy
  wpływ. Contamination, $\nu$ i $\varepsilon$ sterują kompromisem
  Precision--Recall, a F1 jest największe, gdy odsetek oznaczonych punktów
  jest bliski udziałowi anomalii.
  \item Przy niezbalansowaniu Accuracy jest myląca (trywialny klasyfikator
  bywa lepszy od najlepszego detektora), a ROC-AUC optymistyczna; podstawą
  oceny powinny być PR-AUC i F1.
  \item Część trafień jest unikatowa dla jednej metody, a fałszywe alarmy są
  w dużej mierze indywidualne -- to przesłanka do budowy zespołu detektorów
  (zadanie 3).
  \item Jakość detekcji bardziej zależy od zbioru (tego, czy anomalie są
  wyjątkami w sensie rozkładu) niż od wyboru algorytmu: rozrzut PR-AUC
  między zbiorami (0,11--0,99) jest kilkukrotnie większy niż między
  najlepszymi metodami na tym samym zbiorze.
\end{enumerate}

\section{Odtworzenie wyników}
\label{sec:repro}

\begin{verbatim}
python -m task1_classical.eda               # tabela 1, rys. 1-3
python -m task1_classical.run_grid          # siatka -> results/task1/grid.csv
python -m task1_classical.run_best          # najlepsze konfiguracje
python -m task1_classical.decision_regions  # rys. obszarów decyzyjnych
python -m task1_classical.analyze           # pozostałe tabele i rysunki
\end{verbatim}

Zbiory ADBench i DAMI są pobierane automatycznie do \code{data/raw/};
plik \code{arrhythmia.mat} z ODDS (używany w zadaniach 2 i 3) trzeba
pobrać ręcznie, bo serwer ODDS nie pozwala na pobieranie automatyczne.

\appendix
\section{Wyniki wszystkich eksperymentów}
\label{app:grid}

Tabele zawierają średnie po podziałach (część testowa) dla każdej
konfiguracji. F1 dotyczy wbudowanej reguły decyzyjnej metody
(contamination $c$, granica OCSVM lub promień $\varepsilon$ DBSCAN).

{\scriptsize
\setlength{\tabcolsep}{4pt}
\subsection*{Isolation Forest}\tab{grid_IForest}
\subsection*{LOF}\tab{grid_LOF}
\subsection*{COF}\tab{grid_COF}
\subsection*{One-Class SVM}\tab{grid_OCSVM}
\subsection*{DBSCAN}\tab{grid_DBSCAN}
}

% Sources of the benchmark datasets, shared by the three reports.
\begin{thebibliography}{99}
\small

\bibitem{adbench} S.~Han, X.~Hu, H.~Huang, M.~Jiang, Y.~Zhao.
\emph{ADBench: Anomaly Detection Benchmark}. Advances in Neural Information
Processing Systems 35 (NeurIPS 2022), Datasets and Benchmarks Track, 2022.
Repozytorium \url{https://github.com/Minqi824/ADBench},
commit \code{3dac822}.

\bibitem{odds} S.~Rayana. \emph{ODDS Library}. Stony Brook University,
Department of Computer Science, 2016. \url{https://odds.cs.stonybrook.edu}.

\bibitem{dami} G.~O.~Campos, A.~Zimek, J.~Sander, R.~J.~G.~B.~Campello,
B.~Micenková, E.~Schubert, I.~Assent, M.~E.~Houle. \emph{On the evaluation of
unsupervised outlier detection: measures, datasets, and an empirical study}.
Data Mining and Knowledge Discovery 30(4), 891--927, 2016.
\url{https://www.dbs.ifi.lmu.de/research/outlier-evaluation/}.

\bibitem{uci} M.~Kelly, R.~Longjohn, K.~Nottingham. \emph{The UCI Machine
Learning Repository}. University of California, Irvine.
\url{https://archive.ics.uci.edu}.

\bibitem{breastw} W.~H.~Wolberg, O.~L.~Mangasarian. \emph{Multisurface method
of pattern separation for medical diagnosis applied to breast cytology}.
Proceedings of the National Academy of Sciences 87(23), 9193--9196, 1990.
(UCI: Breast Cancer Wisconsin (Original)).

\bibitem{wdbc} W.~N.~Street, W.~H.~Wolberg, O.~L.~Mangasarian. \emph{Nuclear
feature extraction for breast tumor diagnosis}. IS\&T/SPIE Symposium on
Electronic Imaging: Science and Technology, vol.~1905, 861--870, 1993.
(UCI: Breast Cancer Wisconsin (Diagnostic)).

\bibitem{mammography} K.~S.~Woods, C.~C.~Doss, K.~W.~Bowyer, J.~L.~Solka,
C.~E.~Priebe, W.~P.~Kegelmeyer. \emph{Comparative evaluation of pattern
recognition techniques for detection of microcalcifications in mammography}.
International Journal of Pattern Recognition and Artificial Intelligence
7(6), 1417--1436, 1993.

\bibitem{wine} S.~Aeberhard, M.~Forina. \emph{Wine}. UCI Machine Learning
Repository, 1991.

\bibitem{cardio} D.~Ayres-de Campos, J.~Bernardes, A.~Garrido,
J.~Marques-de-Sá, L.~Pereira-Leite. \emph{SisPorto 2.0: a program for
automated analysis of cardiotocograms}. Journal of Maternal-Fetal Medicine
9(5), 311--318, 2000. (UCI: Cardiotocography).

\bibitem{campaign} S.~Moro, P.~Cortez, P.~Rita. \emph{A data-driven approach
to predict the success of bank telemarketing}. Decision Support Systems 62,
22--31, 2014. (UCI: Bank Marketing).

\bibitem{arrhythmia} H.~A.~Güvenir, B.~Acar, G.~Demiröz, A.~Çekin.
\emph{A supervised machine learning algorithm for arrhythmia analysis}.
Computers in Cardiology 1997, 433--436. (UCI: Arrhythmia).

\bibitem{heart} \emph{Statlog (Heart)}. UCI Machine Learning Repository.

\end{thebibliography}


\end{document}
